%% kapa: exam=1 round=23 year=2012 period=2 form=A subject=accounting subject_ko=회계학 seq=29 number=29
%% source: Q-Net 감정평가사 기출문제 (cst00309) artlSeq=5052580 file=raw/1008/5052580/2012년 제23회 감정평가사 제1차시험 기출문제(2교시 A형).pdf
%% conversion: kordoc 4.24.51 -> tools/gen1.py (auto); review=proofread-sonnet; flags=none
\begin{question}{29}
(주)서울은 20×1년 초에 토지를 ￦10,000에 취득하여 사용하기 시작하였다. 20×1년 말과 20×2년 말 현재 동 토지의 공정가치는 각각 ￦11,000과 ￦9,500이다. (주)서울은 매기간 말에 토지를 공정가치로 평가하는 재평가모형을 적용한다. (주)서울의 토지의 회계처리에 관한 설명으로 옳지 않은 것은?
\begin{choices}
  \item 토지재평가가 20×1년과 20×2년 당기손익에 미치는 영향은 없다.
  \item 20×1년 말과 20×2년 말 재평가잉여금의 잔액은 각각 ￦1,000과 ￦0이다.
  \item 만약 (주)서울이 20×2년 초에 동 토지를 모두 처분하였다면, 20×1년 말에 인식한 재평가잉여금은 당기손익으로 재분류할 수 없다.
  \item 만약 (주)서울이 20×3년 초에 동 토지를 ￦9,600에 처분하였다면, (주)서울은 처분이익 ￦100을 당기손익으로 인식한다.
  \item (주)서울의 20×2년 당기손익으로 보고할 토지재평가손실은 ￦500이다.
\end{choices}
\end{question}
\answer{1}
